% 96-12(인쇄 96쪽) — 추첨권의 상금별 장수 표. 원문은 표의 오른쪽 아래 한 칸(상금 0원의 장수)이 찢어져 보이지 않는 그림이라
% tabular 로는 그 찢긴 자리를 나타낼 수 없어 TikZ 로 다시 그렸다. 칸 값은 원문 그대로이고 찢긴 칸은 원문처럼 비워 둔다
% (전체 장수는 문제에서 구하는 값이므로 그림에 적지 않는다).
\begin{tikzpicture}[line width=0.5pt, font=\small]
  \def\cA{2.5}
  \def\cB{4.1}
  % 가로줄 — 머리줄과 값이 보이는 세 행
  \foreach \y in {0,-0.72,-1.36,-2.00} \draw (0,\y) -- (\cB,\y);
  % 마지막 행(상금 0원)은 찢긴 자리까지만
  \draw (0,-2.64) -- (\cA,-2.64);
  \draw (0,-3.28) -- (2.50,-3.28);
  % 세로줄
  \draw (0,0) -- (0,-3.28);
  \draw (\cA,0) -- (\cA,-2.60);
  \draw (\cB,0) -- (\cB,-2.52);
  % 칸 값
  \node at (1.25,-0.36) {\bfseries 상금(원)};
  \node at (3.30,-0.36) {\bfseries 장수};
  \node at (1.25,-1.04) {$1000000$};
  \node at (3.30,-1.04) {$1$};
  \node at (1.25,-1.68) {$100000$};
  \node at (3.30,-1.68) {$10$};
  \node at (1.25,-2.32) {$10000$};
  \node at (3.30,-2.32) {$100$};
  \node at (1.25,-2.96) {$0$};
  % 찢어진 자리 — 오른쪽 아래 한 칸의 가장자리
  \draw plot[smooth] coordinates {(\cA,-2.60) (2.78,-2.52) (3.06,-2.61) (3.36,-2.50) (3.70,-2.59) (\cB,-2.52)};
  \draw plot[smooth] coordinates {(\cA,-2.60) (2.64,-2.78) (2.46,-2.98) (2.62,-3.14) (2.50,-3.28)};
\end{tikzpicture}
